\chapter{Discussion}
\section{Analyse des problèmes rencontrés}
Au cours du projet, deux problèmes ont été découverts et résolus :
\begin{itemize}
    \item \textbf{Problème 1 — Le scan "avant" a été effectué sans le pare-feu actif} : Les règles du pare-feu ont été purgées pour capturer l'état "avant", mais le deuxième scan a été exécuté par erreur avant de réappliquer les règles. La réapplication du script et la réexécution du scan ont résolu le problème.
    \item \textbf{Problème 2 — La règle contre les inondations SYN était trop large} : La règle originale acceptait tous les paquets TCP SYN en dessous de 25/sec, quel que soit le port de destination. La restriction de la règle aux seuls ports 22, 80 et 443 a fait apparaître les ports bloqués comme "filtrés" au lieu de "fermés".
\end{itemize}

\section{ufw — L'alternative plus simple}
\texttt{ufw} (Uncomplicated Firewall) est un outil de plus haut niveau qui écrit automatiquement les règles \texttt{iptables}.
\begin{lstlisting}[language=bash]
sudo ufw default deny incoming
sudo ufw default allow outgoing
sudo ufw allow ssh
sudo ufw allow http
sudo ufw allow https
sudo ufw limit ssh
sudo ufw enable
\end{lstlisting}
\texttt{ufw} est plus facile à utiliser mais moins flexible que l'écriture directe de règles \texttt{iptables}. Travailler avec \texttt{iptables} brutes offre une compréhension plus approfondie du filtrage de paquets.
\lipsum[9-10]
